% ---------------------------------------------------------------------
%  CHƯƠNG 1 — MỞ ĐẦU
% ---------------------------------------------------------------------
\chapter{Mở đầu}
\label{ch:modau}

\begin{muctieu}
Chương mở đầu trình bày bối cảnh, lý do lựa chọn đề tài, mục tiêu, phạm vi
và đối tượng thực hiện của đồ án. Nội dung chương nhằm làm rõ nhu cầu xây
dựng một trợ lý AI có khả năng cá nhân hoá để hỗ trợ người khiếm thị trong
quá trình tiếp cận thông tin, nhận diện không gian/vật thể và tương tác với hệ thống.
\end{muctieu}

\section{Lý do chọn đề tài}

Trong quá trình phát triển của trí tuệ nhân tạo, các trợ lý AI ngày càng
được ứng dụng rộng rãi để hỗ trợ con người trong việc tìm kiếm thông tin,
giao tiếp và thực hiện các tác vụ hằng ngày. Tuy nhiên, phần lớn các trợ lý
AI hiện nay được thiết kế theo hướng phục vụ người dùng nói chung, chủ yếu dựa
trên giao diện trực quan và hình ảnh, chưa tập trung đầy đủ vào nhu cầu tiếp
cận thông tin và tương tác bằng âm thanh/giọng nói của người khiếm thị.

Người khiếm thị gặp rất nhiều khó khăn trong các tình huống giao tiếp và sinh
hoạt hằng ngày vốn phụ thuộc vào thị giác, chẳng hạn như đọc tài liệu văn bản,
nhận diện các vật thể xung quanh, quan sát biển hiệu hay di chuyển trong môi
trường lạ. Vì vậy, một hệ thống hỗ trợ dựa trên trí tuệ nhân tạo có khả năng
chuyển đổi thông tin thị giác thành âm thanh/lời nói hoặc phản hồi thông minh
qua giọng nói có thể giúp giảm bớt những trở ngại trong quá trình tiếp cận công nghệ
và hòa nhập cộng đồng~\cite{WHO2019WorldVision,WHO2022AssistiveTech}.

Bên cạnh khả năng cung cấp và hỗ trợ xử lý thông tin, khả năng điều chỉnh cách
tương tác cũng có ý nghĩa quan trọng. Mỗi người dùng khiếm thị có thể có nhu cầu
khác nhau về cỡ chữ, thao tác, giọng đọc và mức độ chi tiết. Trong phạm vi
nguyên mẫu này, các lựa chọn được áp dụng trong phiên làm việc; hệ thống chưa
lưu hồ sơ cá nhân lâu dài.

Xuất phát từ những vấn đề trên, đồ án lựa chọn đề tài ``Xây dựng trợ lý AI
cá nhân hoá hỗ trợ người khiếm thị'' nhằm nghiên cứu khả năng ứng dụng
trí tuệ nhân tạo vào một bài toán có tính thực tiễn cao, hướng đến việc gia tăng
khả năng tiếp cận công nghệ và nâng cao chất lượng cuộc sống cho người dùng có
nhu cầu đặc thù.

\section{Mục tiêu của đồ án}

Đồ án hướng tới các mục tiêu sau:
\begin{itemize}[leftmargin=1.6em, itemsep=2pt]
  \item Phân tích nhu cầu và các khó khăn trong quá trình tiếp cận thông tin,
        quan sát và tương tác với trợ lý AI của người khiếm thị.
        
  \item Xây dựng một trợ lý AI có khả năng tiếp nhận yêu cầu bằng giọng nói
        hoặc dữ liệu đầu vào (văn bản, hình ảnh), xử lý nội dung và đưa ra
        phản hồi bằng âm thanh/giọng nói phù hợp.
        
  \item Bổ sung tuỳ chọn tương tác trong phiên như chế độ trợ năng, giọng đọc,
        mức chi tiết và ngữ cảnh hội thoại; không đặt mục tiêu xây dựng hồ sơ
        người dùng lâu dài trong phiên bản này.
        
  \item Thiết kế giao diện dễ tiếp cận với chế độ chữ/nút lớn, focus bàn phím,
        robot hướng dẫn và khả năng nghe kết quả bằng giọng đọc.
        
  \item Kiểm thử các chức năng phần mềm đã xây dựng; xác định dữ liệu và phương pháp
        cần có để đánh giá riêng độ chính xác, độ trễ và mức độ đáp ứng bài toán.
\end{itemize}

\section{Phạm vi và đối tượng}

Đồ án tập trung xây dựng một nguyên mẫu trợ lý thị giác AI hướng đến hỗ trợ
người khiếm thị tiếp nhận thông tin trong ảnh, đọc văn bản và tương tác với
trợ lý bằng tiếng Việt. Đối tượng hướng đến là người mù hoặc người có thị lực
kém cần công cụ hỗ trợ đọc và nhận biết thông tin môi trường.

Phạm vi của đồ án tập trung vào nguyên mẫu Streamlit gồm phân tích ảnh Gemini,
OCR cục bộ bằng Tesseract, hỏi đáp theo ảnh, lệnh thoại, đọc phản hồi và các
tuỳ chọn trợ năng. STT sử dụng Google Web Speech qua Internet; TTS ưu tiên
Edge TTS và có Windows SAPI5 dự phòng. Các tuỳ chọn cá nhân hoá chỉ tồn tại
trong phiên làm việc, không gồm tài khoản hoặc hồ sơ lâu dài.

Đồ án không hướng tới việc thay thế hoàn toàn các công cụ hỗ trợ di chuyển chuyên dụng
(như gậy trắng, chó dẫn đường) hay các thiết bị y tế chuyên khoa, không giải quyết
toàn bộ các vấn đề y tế liên quan đến thị giác và chưa đặt mục tiêu triển khai hệ thống
ở quy mô thương mại/sản xuất lớn. Các chức năng và phạm vi hỗ trợ cụ thể được giới hạn
theo thời gian và nguồn lực thực hiện đồ án.

\begin{luuy}
Phạm vi của đồ án được giới hạn ở việc xây dựng và đánh giá một nguyên mẫu
trợ lý AI cá nhân hoá hỗ trợ người khiếm thị. Việc giới hạn phạm vi giúp
tập trung vào các chức năng cốt lõi, bảo đảm hệ thống có thể được triển khai,
kiểm thử và đánh giá trong thời gian thực hiện đồ án.
\end{luuy}

\section{Bố cục báo cáo}

Phần còn lại của báo cáo được tổ chức như sau. Chương~\ref{ch:coso} trình
bày cơ sở lý thuyết và các công nghệ được sử dụng trong quá trình xây dựng
trợ lý AI. Chương~\ref{ch:thietke} phân tích bài toán, yêu cầu hệ thống và
thiết kế giải pháp. Chương~\ref{ch:caidat} mô tả quá trình cài đặt và xây
dựng các chức năng của hệ thống. Chương~\ref{ch:ketqua} trình bày kết quả
thực nghiệm và đánh giá hệ thống. Chương~\ref{ch:ketluan} tổng kết những
kết quả đạt được, các hạn chế và hướng phát triển của đề tài.